% !TeX program = xelatex
% UTF-8. Compile with XeLaTeX from this directory or the repository root.
% Source: Google Drive / Notebook / maternal diseases that affect neonates.
% pic01.jpg = 20261016_105619571.jpg; Drive ID 1Ecofn5uG9M89cQ7bsS_xKzjjNsvyVuLz.
% pic02.jpg = 20261016_105626746.jpg; Drive ID 1H7O9FWIUAOx1T6n33I0YIxBwrSXoYLSD.
% Define \MaternalDiseasesInNotebook before inputting into a XePersian document.
\ifdefined\MaternalDiseasesInNotebook
  \let\MDfinish\relax
\else
  \def\MDfinish{\end{document}}
  \documentclass[12pt,a4paper]{article}
  \usepackage[margin=22mm]{geometry}
  \usepackage{graphicx}
  \usepackage{enumitem}
  \usepackage{xepersian}
  \IfFontExistsTF{Amiri}{\settextfont{Amiri}}{\settextfont{DejaVu Sans}}
  \setlatintextfont{DejaVu Serif}
  \setlength{\parindent}{0pt}
  \setlength{\parskip}{5pt}
  \setlength{\emergencystretch}{3em}
  \begin{document}
\fi
\providecommand{\MDmark}[1]{\textbf{#1}}
\providecommand{\MDunclear}[1]{\textbf{[خوانش نامطمئن: #1]}}

\begin{center}
  {\Large\bfseries ۵۹ ـ بیماری‌های مادری مؤثر بر نوزاد}\par
  \lr{Maternal diseases that affect neonates}
\end{center}

\textbf{یادداشت رونویسی:}
متن فارسی و انگلیسی و شاخه‌های دست‌نویس به ترتیب دو تصویر آمده‌اند؛
عبارت‌های مشخص‌شده با ماژیک پررنگ شده‌اند و خوانش نامطمئن در متن مشخص است.
تصاویر کامل منبع در پایان همین سند آمده‌اند.
سربرگ چاپی هر دو صفحه: \lr{Subject:}، \lr{Year:}، \lr{Month:}، \lr{Date:}؛
خانه‌های تاریخ خالی‌اند. موضوع دست‌نویس صفحهٔ اول:
«۵۹ ـ بیماری‌های مادری مؤثر بر نوزاد»؛ موضوع دست‌نویس صفحهٔ دوم: «۸۴».
اعداد چاپی کنار سطرهای هر دو برگه: ۱ تا ۹؛ شماره‌های دست‌نویس پایین صفحات: ۱ و ۲.

\section*{تصویر ۱ ـ صفحهٔ ۱}
\subsection*{سندرم آنتی‌فسفولیپید (\lr{APS})}
\begin{itemize}
  \item \MDmark{ترومبوفیلی + سقط مکرر}.
  \item اثر پروترومبوتیک آنتی‌بادی‌ها بر عملکرد جفت.
  \item واسکولوپاتی، انفارکت و ترومبوز در مادر.
  \item \MDmark{اختلال رشد جنین} و \MDmark{نارسایی جفت}، پره‌اکلامپسی، زایمان زودرس.
  \item ترومبوز در جنین.
\end{itemize}

\subsection*{\lr{ITP}}
\begin{itemize}
  \item در ۲ مورد از هر ۱۰۰۰ تولد؛ \MDmark{\lr{IgG}}.
  \item \MDmark{ترومبوسایتوپنی در نوزاد}.
  \item ریسک \lr{ICH}.
  \item در صورت خونریزی: \MDmark{پلاکت اشعه‌دیدهٔ \lr{Single donor}} دریافت کند.
  \item \lr{IVIG} برای نوزاد ممکن است کمک‌کننده باشد.
  \item معمولاً ترومبوسایتوپنی در نوزاد ظرف \MDmark{۴ تا ۶ هفته} رفع می‌شود.
\end{itemize}

\subsection*{\lr{SLE}}
\begin{itemize}
  \item \MDmark{\lr{IgG}}؛ \lr{anti-Ro (SS-A) Ab}، \lr{anti-La (SS-B) Ab}.
  \item آنتی‌بادی‌ها از جفت عبور می‌کنند و به جنین آسیب می‌زنند.
  \item جدی‌ترین عارضه: \MDmark{\lr{Congenital Heart Block}}؛
        \MDmark{معمولاً درجهٔ ۳} (کامل).
  \item اغلب نوزادان نیاز به \MDmark{\lr{pacing}} پیدا می‌کنند.
\end{itemize}
\subsubsection*{سایر علائم}
\begin{itemize}
  \item ضایعات پوستی: پلاک‌های اریتماتو.
  \item آتروفی عضلات \MDunclear{نترال؛ واژهٔ پس از «عضلات»}.
  \item \lr{Scaling} خفیف: چشم‌ها، صورت، اسکالپ.
  \item \MDmark{ترومبوسایتوپنی}.
  \item \MDmark{همولیز خودایمن}.
  \item درگیری کبدی.
\end{itemize}

\subsection*{هایپرتیروئیدی نوزادی}
\begin{itemize}
  \item ۰٫۴\% زنان باردار هایپرتیروئیدی کلینیکال دارند.
  \item آنتی‌بادی‌های تحریک‌کنندهٔ تیروئید از جفت عبور می‌کنند.
  \item علائم ظرف \MDmark{۱۲ تا ۴۸ ساعت اول تولد} نمایان می‌شود.
  \item \lr{IUGR}، نارس بودن، گواتر (انسداد تراشه)، اگزوفتالمی،
        کرانیوسینوستوزیس، هایپوگلیسمی، ترومبوسایتوپنی، ارگانومگالی.
  \item ظرف \MDmark{۲ تا ۴ ماه} رفع می‌شود.
  \item درمان: \MDmark{پروپیل تیواوراسیل} ـ قطرهٔ ید ـ پروپرانولول.
\end{itemize}

\clearpage
\section*{تصویر ۲ ـ صفحهٔ ۲}
\subsection*{\lr{DM}}
\begin{itemize}
  \item دیابت کنترل‌نشدهٔ مادر: هایپرپلازی جزایر لانگرهانس جنین.
  \item هایپرانسولینمی جنینی: ماکروزومی؛ هایپوگلیسمی ناشتا نوزاد.
\end{itemize}

\subsection*{داروها}
\begin{itemize}
  \item \MDmark{استامینوفن در بارداری ایمن است.}
  \item \lr{NSAIDs}: خطر بسته شدن زودهنگام \lr{DA}؛ \MDmark{هایپرتنشن ریوی}.
  \item ایبوپروفن: \MDmark{کاهش حجم مایع آمنیوتیک}.
  \item کاربامازپین (ضدتشنج): \MDmark{نقص لولهٔ عصبی} (\lr{NTD}).
\end{itemize}

\textbf{حاشیهٔ تصویر ۲:}
قطعهٔ کوچک نوشتهٔ لاتین در محل سوراخ سمت چپ برگه: \lr{NS}
(بخش باقی‌ماندهٔ نوشته؛ ادامهٔ آن در تصویر دیده نمی‌شود).

\clearpage
\section*{تصاویر کامل منبع}
\IfFileExists{images/pic01.jpg}{\def\MDimagebase{images/}}{%
  \def\MDimagebase{pediatrics/ped-neonatology/maternal diseases that affect neonates/images/}}
\begin{center}
  \textbf{تصویر ۱ ـ \lr{pic01.jpg}}\par
  \includegraphics[width=\linewidth,height=0.87\textheight,keepaspectratio]{\MDimagebase pic01.jpg}
\end{center}
\clearpage
\begin{center}
  \textbf{تصویر ۲ ـ \lr{pic02.jpg}}\par
  \includegraphics[width=\linewidth,height=0.87\textheight,keepaspectratio]{\MDimagebase pic02.jpg}
\end{center}
\MDfinish
